% !PRINTED 279
\weberpage{279}
%<<CONT
to be set, from which, by forming $R\,A\,R^{-1}$, the generating substitutions of the octahedral group result in real form:
\[
\Theta = \begin{pmatrix} 2, & 1 \\ 3, & 2 \end{pmatrix}, \
\psi = \begin{pmatrix} 2, & 2 \\ 1, & -2 \end{pmatrix}, \
\omega = \psi\,\Theta = \begin{pmatrix} 3, & -1 \\ 3, & -3 \end{pmatrix}
\]
\weq{11}{
\chi = \begin{pmatrix} 0, & 2 \\ 3, & 1 \end{pmatrix}, \
\chi^2 = \begin{pmatrix} -1, & 2 \\ 3, & 0 \end{pmatrix}.
}
It is now easy, following \S.\,57, to form the complete group, which it is not necessary to carry out further.

\gesperrt{The octahedral group for $p=7$ is a divisor of $L_7$ of index 7.}

V. \gesperrt{Icosahedral group.}

In the case $p=5$, $L_p$ is itself an icosahedral group. For other values of $p$, an icosahedral group can be contained in $L_p$ only if $p^2-1$ is divisible by 5, that is, $p\equiv\pm 1\ (\mathrm{mod}\ 5)$. Then from \S.\,58 we obtain the generating substitutions of such a group.

If, for brevity, we set
\[
\gamma^{\frac{p^2-1}{10}} = \varrho,
\]
then we have to substitute $\varrho^2$ for $\varepsilon$ in the formulae of \S.\,58, and obtain
\weq{12}{
\Theta = \begin{pmatrix} \varrho, & 0 \\ 0, & \varrho^{-1} \end{pmatrix}, \
\psi = \begin{pmatrix} 0, & 1 \\ -1, & 0 \end{pmatrix}, \quad
\chi = \begin{pmatrix} \dfrac{-1}{\varrho^2-\varrho^{-2}}, & \dfrac{1}{\varrho-\varrho^{-1}} \\[2mm] \dfrac{1}{\varrho-\varrho^{-1}}, & \dfrac{1}{\varrho^2-\varrho^{-2}} \end{pmatrix}\ [\S.\,58,\,(26)].
}
If $p\equiv 1\ (\mathrm{mod}\ 5)$, then $\varrho$, and with it the whole group, is real. But if $p\equiv -1\ (\mathrm{mod}\ 5)$, then $\varrho$ and $\varrho^{-1}$ are conjugate, and consequently the group found is contained in $\Gamma_p$, and in order to obtain the icosahedral group in $L_p$ we must first transform according to \S.\,69.

\gesperrt{For $p=11$ we obtain immediately a real icosahedral group of index 11.}
